%% Front matter shared by both versions.
%% The abstract below is the journal version. The plain-text ASCII abstract
%% for the arXiv submission form lives in docs/ARXIV-METADATA.txt and is
%% pasted into the form by hand: arXiv does not read it from the PDF.

\title[Hyperbolicity of renormalization at a non-integer exponent]
      {Existence and hyperbolicity of the period-doubling renormalization
       fixed point for non-integer critical exponents}

\ifx\ANONYMOUS\undefined
  \author{Yevhen Borzenkov}
  \address{Independent researcher, Bratislava, Slovakia}
  \urladdr{https://orcid.org/0009-0002-7946-589X}
\fi

\subjclass[2020]{Primary 37E20; Secondary 37E05, 37F25, 47H10, 65G20, 65G30, 37-04}
\keywords{period-doubling renormalization; Feigenbaum--Cvitanovi\'{c} equation;
unimodal maps; non-integer critical exponent; computer-assisted proof;
interval arithmetic; Newton--Kantorovich theorem; Rouch\'{e}'s theorem;
hyperbolic fixed point; Feigenbaum constants}

\begin{abstract}
We prove, by a computer-assisted argument with validated interval arithmetic,
that for every critical exponent $\ell \in [2,4]$ the Feigenbaum--Cvitanovi\'{c}
equation has a solution $\Gs(\ell)$, analytic on an explicit disk and normalized
by $\Gs(0) = 1$, and that the derivative of the period-doubling renormalization
at $\Gs(\ell)$, restricted to the normalization slice $G(0) = 1$, has exactly one
eigenvalue outside the unit circle; this
eigenvalue $\delta(\ell)$ is real and simple, and both $\Gs(\ell)$ and
$\delta(\ell)$ depend real-analytically on $\ell$. The solution lies in a ball of
radius $8.49 \cdot 10^{-13}$ around a computable parametrized center and is
unique on the normalization slice within $1.75 \cdot 10^{-7}$. The proof verifies
a Newton--Kantorovich gate over 46 intervals covering $[2,4]$ and an
operator-valued Rouch\'{e} inequality over 706 intervals, with resolvent bound
$3.7961$ on the unit circle and worst acceptance value $0.9919 < 1$. Unlike the
computer-free proofs of Gorbovickis and Yampolsky and of Cruz and Smania, which
hold in neighborhoods of even integers of unspecified width, the covering is
explicit and uniform in $\ell$ and includes interior points such as
$\ell = 5/2$ and $\ell = 3$. Independence of the spectrum from the choice of the
layer $(c, r)$ is proved for an explicitly certified family of weighted
$\ell^{1}$ spaces: a rectangle of layers around the working one, fourteen
further layers, and a corridor joining the rectangle to one of the farthest of them.
\end{abstract}

\maketitle
